\hypertarget{software-requirements-specification-codeguard-ai}{%
\chapter{Software Requirements Specification --- CodeGuard
AI}\label{software-requirements-specification-codeguard-ai}}

\textbf{Project:} Code Plagiarism Detector Using AST and Behavioural
Fingerprinting \textbf{Institution:} Invertis University, Bareilly \textperiodcentered{} AY
2026--2027 \textbf{Document version:} 1.0, 03 November 2026
\textbf{Intended placement:} Report Appendix A, summarised in Chapter 3

\begin{draftnote}

\textbf{Nature of this document.} This is an \textbf{as-built
specification}: it was reconstructed from the implemented system, its
tests and its configuration, rather than written in advance of
construction. It is presented as such deliberately. Every requirement
below is traceable to code that exists, and no requirement describes
behaviour that has not been built. Where the implementation falls short
of what a requirement would ideally state, the shortfall is recorded in
the requirement itself rather than omitted.

Requirements are identified as FR-n (functional), NFR-n (non-functional)
and DR-n (data). Identifiers are stable; do not renumber them, because
Chapter 6's test table refers to them.

\end{draftnote}

\begin{center}\rule{0.5\linewidth}{0.5pt}\end{center}

\hypertarget{introduction}{%
\section{Introduction}\label{introduction}}

\hypertarget{purpose}{%
\subsection{Purpose}\label{purpose}}

CodeGuard AI ranks programming assignment submissions by the likelihood
that they warrant human review for academic misconduct, and attaches to
each ranking the evidence that produced it. It is an evidence system and
not a verdict machine: it produces an ordering and a justification, and
a member of faculty makes the decision.

This document specifies what the system does, for whom, under what
constraints, and with what limits.

\hypertarget{scope}{%
\subsection{Scope}\label{scope}}

The system addresses three forms of misconduct in submitted source code:

\begin{enumerate}
\def\labelenumi{\arabic{enumi}.}
\tightlist
\item
  \textbf{Collusion} --- two or more students submitting structurally
  equivalent work. Addressed by Layer 1, structural AST similarity.
\item
  \textbf{Outsourcing} --- a student submitting work written by a third
  party, with no second copy in the cohort. Addressed in design by Layer
  2, behavioural fingerprinting. \textbf{See NFR-9 for the measured
  limitation on this claim.}
\item
  \textbf{Machine-generated submission} --- work produced by a language
  model. \textbf{Out of scope in this version}; see \S{}2.6.
\end{enumerate}

Supported languages are Python and Java. The system handles single-file
submissions.

\hypertarget{definitions}{%
\subsection{Definitions}\label{definitions}}

\begin{longtable}[]{@{}
  >{\raggedright\arraybackslash}p{(\columnwidth - 2\tabcolsep) * \real{0.5000}}
  >{\raggedright\arraybackslash}p{(\columnwidth - 2\tabcolsep) * \real{0.5000}}@{}}
\toprule\noalign{}
\begin{minipage}[b]{\linewidth}\raggedright
term
\end{minipage} & \begin{minipage}[b]{\linewidth}\raggedright
meaning
\end{minipage} \\
\midrule\noalign{}
\endhead
\bottomrule\noalign{}
\endlastfoot
\textbf{RPS} & Review Priority Score. The composed, renormalised score a
submission is ranked by. \\
\textbf{Layer 1} & Structural detection: AST comparison between a
submission and cohort candidates. \\
\textbf{Layer 2} & Behavioural detection: comparison of a submission's
style against its own author's baseline. \\
\textbf{Layer 3} & Detection of machine-generated code. Not implemented;
see \S{}2.6. \\
\textbf{Anchor} & A past submission trustworthy enough to describe a
student's style. \\
\textbf{Baseline} & A per-student, per-language style profile built from
anchors. \\
\textbf{Cohort} & The other students whose work provides comparison
candidates and distributional statistics. \\
\textbf{Provenance} & Whether an assignment was invigilated or taken
home, which determines anchor trust. \\
\end{longtable}

\hypertarget{references}{%
\subsection{References}\label{references}}

\begin{itemize}
\tightlist
\item
  \texttt{docs/PROJECT\_STATUS.md} --- the as-built inventory this
  document derives from
\item
  \texttt{docs/PROJECT\_NOTES.md} --- the decision record, including
  retractions
\item
  \texttt{docs/TESTING.md} --- the suite contract
\item
  \texttt{apps/detector/measurements/README.md} --- the measurement
  record and its limits
\item
  \texttt{docs/report/references.bib} --- academic citations
\end{itemize}

\begin{center}\rule{0.5\linewidth}{0.5pt}\end{center}

\hypertarget{overall-description}{%
\section{Overall description}\label{overall-description}}

\hypertarget{product-perspective}{%
\subsection{Product perspective}\label{product-perspective}}

Three tiers and one detection service:

\begin{itemize}
\tightlist
\item
  \textbf{Web client} (\texttt{apps/web}) --- React and Vite. Holds no
  credentials and talks only to the API.
\item
  \textbf{API} (\texttt{apps/api}) --- Node, Express, TypeScript,
  Mongoose. Holds all credentials and is the only component that reaches
  a data store.
\item
  \textbf{Detector} (\texttt{apps/detector}) --- Python, FastAPI,
  tree-sitter, APTED. Stateless; parses and compares source, and
  extracts style features.
\item
  \textbf{Infrastructure} --- MongoDB (documents), Redis (job queue),
  MinIO-compatible object store (submitted files).
\item
  \textbf{Worker} --- the same codebase as the API started on a
  different entry point, so a queued job and a request share one set of
  services rather than two implementations that can diverge.
\end{itemize}

\hypertarget{user-classes}{%
\subsection{User classes}\label{user-classes}}

\begin{longtable}[]{@{}
  >{\raggedright\arraybackslash}p{(\columnwidth - 2\tabcolsep) * \real{0.5000}}
  >{\raggedright\arraybackslash}p{(\columnwidth - 2\tabcolsep) * \real{0.5000}}@{}}
\toprule\noalign{}
\begin{minipage}[b]{\linewidth}\raggedright
class
\end{minipage} & \begin{minipage}[b]{\linewidth}\raggedright
may
\end{minipage} \\
\midrule\noalign{}
\endhead
\bottomrule\noalign{}
\endlastfoot
\textbf{Student} & submit work, view their own submission status and any
recorded decision, answer a flag, nominate past work as an anchor \\
\textbf{Faculty} & create courses and assignments, enrol students, view
the ranked review queue and full evidence, record decisions, accept
nominations, re-run detection \\
\textbf{Admin} & all of the above, plus create accounts with roles and
edit detection configuration \\
\end{longtable}

A student never sees a score. This is a design decision and not an
oversight: a ranking a student can read is a ranking a student can tune
against.

\hypertarget{operating-environment}{%
\subsection{Operating environment}\label{operating-environment}}

Developed and run on Windows with Docker Desktop providing MongoDB,
Redis and the object store via \texttt{infra/docker-compose.yml}. CI
runs on Ubuntu with the same services. The detector targets Python 3.11.

\hypertarget{design-and-implementation-constraints}{%
\subsection{Design and implementation
constraints}\label{design-and-implementation-constraints}}

\begin{itemize}
\tightlist
\item
  Single-file submissions only; no archive or multi-file project
  support.
\item
  Two languages only, enforced by extension allow-list.
\item
  The detector accepts at most 50 comparison candidates per request, and
  compares at most 10 of them in full (\S{}3.3, NFR-3).
\item
  The API is the only holder of credentials; no credential appears in
  the web client or the detector.
\end{itemize}

\hypertarget{assumptions-and-dependencies}{%
\subsection{Assumptions and
dependencies}\label{assumptions-and-dependencies}}

\begin{itemize}
\tightlist
\item
  A student's past work is available and attributable to them, which is
  what makes a baseline possible. Where it is not, Layer 2 declines to
  score.
\item
  Invigilated work is genuinely the student's own. The anchor provenance
  model rests on this and cannot verify it.
\item
  Cohort members are solving comparable problems, so that distributional
  statistics are meaningful.
\end{itemize}

\hypertarget{explicitly-excluded-from-this-version}{%
\subsection{Explicitly excluded from this
version}\label{explicitly-excluded-from-this-version}}

\textbf{Layer 3, detection of machine-generated code, is not
implemented.} The decision was taken on 06 October 2026 and is recorded
rather than hidden: every detection result stores
\texttt{aiContent:\ \{\ status:\ "skipped",\ reason:\ "Layer\ 3\ is\ not\ implemented\ yet"\ \}}.

The reason is a requirements problem and not a scheduling one. The
published approach this layer would have followed reports that
classifier accuracy depends on training data resembling the target code,
which requires a labelled corpus of student work from the course being
policed. No such corpus exists, and manufacturing one would mean
treating synthetic submissions as ground truth. A layer whose output
cannot be characterised is worse than an absent layer in a system whose
premise is auditable evidence, because a reviewer would have no basis on
which to weigh it.

Also excluded: real-time push notification of queue changes (the queue
is polled), and side-by-side AST diff rendering of two sources.

\begin{center}\rule{0.5\linewidth}{0.5pt}\end{center}

\hypertarget{specific-requirements}{%
\section{Specific requirements}\label{specific-requirements}}

\hypertarget{functional-requirements}{%
\subsection{Functional requirements}\label{functional-requirements}}

\hypertarget{authentication-and-accounts}{%
\subsubsection{Authentication and
accounts}\label{authentication-and-accounts}}

\begin{itemize}
\tightlist
\item
  \textbf{FR-1} The system shall allow public self-registration, which
  shall create an account with the student role and no other.
\item
  \textbf{FR-2} The system shall authenticate by email and password and
  issue a short-lived access token and a longer-lived refresh token.
\item
  \textbf{FR-3} The system shall allow a held refresh token to be
  exchanged for a new access token, and shall allow logout to invalidate
  it.
\item
  \textbf{FR-4} The system shall allow an authenticated user to retrieve
  their own account.
\item
  \textbf{FR-5} The system shall allow an admin, and only an admin, to
  create an account with a chosen role, and to amend an existing
  account.
\item
  \textbf{FR-6} The system shall allow an admin to list accounts, paged,
  filtered by role and active state, and searchable by term.
\end{itemize}

\hypertarget{courses-enrolment-and-configuration}{%
\subsubsection{Courses, enrolment and
configuration}\label{courses-enrolment-and-configuration}}

\begin{itemize}
\tightlist
\item
  \textbf{FR-7} The system shall allow faculty to create a course and to
  amend it.
\item
  \textbf{FR-8} The system shall allow faculty to enrol a student on a
  course and to remove one.
\item
  \textbf{FR-9} The system shall expose, per course, the detection
  configuration in force, and shall allow an authorised user to amend
  it.
\item
  \textbf{FR-10} Detection configuration shall support per-language
  overrides, so that one course may weight Python and Java differently.
\end{itemize}

\hypertarget{assignments}{%
\subsubsection{Assignments}\label{assignments}}

\begin{itemize}
\tightlist
\item
  \textbf{FR-11} The system shall allow faculty to create an assignment
  carrying its language and its provenance (invigilated or take-home),
  and to amend it.
\item
  \textbf{FR-12} The system shall expose, per assignment, the review
  queue ordered by RPS descending.
\end{itemize}

\hypertarget{submission}{%
\subsubsection{Submission}\label{submission}}

\begin{itemize}
\tightlist
\item
  \textbf{FR-13} The system shall accept a single source file per
  submission, rejecting any file whose extension is outside the
  configured allow-list or whose size exceeds the configured maximum.
\item
  \textbf{FR-14} The system shall store the submitted file in the object
  store and retain a content hash of it.
\item
  \textbf{FR-15} The system shall enqueue detection on upload, debounced
  by job identity, so that many uploads arriving at a deadline collapse
  into a bounded number of detection runs rather than one run per
  upload.
\item
  \textbf{FR-16} The system shall allow a student to retrieve their own
  submission and its file, and faculty to retrieve any submission within
  their courses.
\end{itemize}

\hypertarget{detection}{%
\subsubsection{Detection}\label{detection}}

\begin{itemize}
\tightlist
\item
  \textbf{FR-17} The system shall compute a structural similarity score
  for a submission against candidates drawn from the cohort.
\item
  \textbf{FR-18} The system shall compute a behavioural score for a
  submission against that student's own baseline for the submission's
  language, where a baseline of sufficient quality exists.
\item
  \textbf{FR-19} The system shall compose the available layer scores
  into one RPS, \textbf{renormalised over the layers that actually ran}.
  An absent layer shall contribute nothing, rather than contributing a
  zero, because a zero is a finding and an absent layer has made no
  finding.
\item
  \textbf{FR-20} The system shall store, on every result, the weights
  and review threshold that produced it, so that the result remains
  interpretable after the configuration changes.
\item
  \textbf{FR-21} The system shall store both layers' evidence on the
  result, not merely their scores, so that the evidence view is
  reproducible from stored data rather than recomputed on demand.
\item
  \textbf{FR-22} The system shall allow faculty to re-run detection on a
  submission; a re-run shall supersede the previous result by revision
  without destroying it, because a decision may already have been
  recorded against the old one.
\end{itemize}

\hypertarget{baseline-integrity}{%
\subsubsection{Baseline integrity}\label{baseline-integrity}}

\begin{itemize}
\tightlist
\item
  \textbf{FR-23} Only invigilated work, or work a student nominates and
  faculty accept, shall be eligible to describe that student's style.
\item
  \textbf{FR-24} The system shall allow a student to nominate a past
  submission as an anchor and to withdraw that nomination, and faculty
  to accept it.
\item
  \textbf{FR-25} Where fewer than \texttt{minAnchorsForBaseline} (3)
  anchors exist, Layer 2 shall decline to score rather than scoring on a
  guess.
\item
  \textbf{FR-26} A thin or inconsistent baseline shall attenuate Layer
  2's effective weight rather than being trusted at full strength.
\item
  \textbf{FR-27} Where a student's own work varies less than the
  cohort's, the result shall carry a low-variance flag, because
  z-scoring penalises such a student.
\item
  \textbf{FR-28} Where the pooled cohort spread has collapsed, the
  result shall carry a cohort-saturation flag, because every z-score is
  then a real difference divided by an invented floor and the layer has
  stopped discriminating.
\end{itemize}

\hypertarget{review-and-the-students-view}{%
\subsubsection{Review and the student's
view}\label{review-and-the-students-view}}

\begin{itemize}
\tightlist
\item
  \textbf{FR-29} The system shall allow faculty to record a decision
  against a submission.
\item
  \textbf{FR-30} The system shall expose to a student a view of their
  own submission carrying its status and any recorded decision,
  \textbf{and no score}.
\item
  \textbf{FR-31} The system shall allow a student to submit a written
  explanation against a recorded decision.
\end{itemize}

\hypertarget{audit}{%
\subsubsection{Audit}\label{audit}}

\begin{itemize}
\tightlist
\item
  \textbf{FR-32} The system shall record every state change affecting a
  student's record in an audit log held separately from the documents it
  describes, so that the record of who changed what cannot be amended by
  the same operation that changes the thing.
\end{itemize}

\hypertarget{external-interfaces}{%
\subsection{External interfaces}\label{external-interfaces}}

35 HTTP endpoints. The API surface is 32 routes across four routers; the
detector exposes 3.

\begin{longtable}[]{@{}
  >{\raggedright\arraybackslash}p{(\columnwidth - 4\tabcolsep) * \real{0.3333}}
  >{\raggedright\arraybackslash}p{(\columnwidth - 4\tabcolsep) * \real{0.3333}}
  >{\raggedright\arraybackslash}p{(\columnwidth - 4\tabcolsep) * \real{0.3333}}@{}}
\toprule\noalign{}
\begin{minipage}[b]{\linewidth}\raggedright
router
\end{minipage} & \begin{minipage}[b]{\linewidth}\raggedright
count
\end{minipage} & \begin{minipage}[b]{\linewidth}\raggedright
endpoints
\end{minipage} \\
\midrule\noalign{}
\endhead
\bottomrule\noalign{}
\endlastfoot
\texttt{authRoutes.ts} & 8 & \texttt{POST\ /register},
\texttt{POST\ /login}, \texttt{POST\ /refresh}, \texttt{POST\ /logout},
\texttt{GET\ /me}, \texttt{GET\ /users}, \texttt{POST\ /users},
\texttt{PATCH\ /users/:userId} \\
\texttt{courseRoutes.ts} & 8 & \texttt{GET\ /},
\texttt{GET\ /:courseId}, \texttt{POST\ /}, \texttt{PATCH\ /:courseId},
\texttt{POST\ /:courseId/students},
\texttt{POST\ /:courseId/students/remove},
\texttt{GET\ /:courseId/config}, \texttt{PATCH\ /:courseId/config} \\
\texttt{assignmentRoutes.ts} & 5 & \texttt{GET\ /}, \texttt{POST\ /},
\texttt{GET\ /:assignmentId}, \texttt{PATCH\ /:assignmentId},
\texttt{GET\ /:assignmentId/results} \\
\texttt{submissionRoutes.ts} & 11 & \texttt{POST\ /}, \texttt{GET\ /},
\texttt{GET\ /:submissionId}, \texttt{GET\ /:submissionId/file},
\texttt{GET\ /:submissionId/result},
\texttt{GET\ /:submissionId/student-view},
\texttt{PATCH\ /:submissionId/explanation},
\texttt{PATCH\ /:submissionId/review},
\texttt{POST\ /:submissionId/nominate},
\texttt{DELETE\ /:submissionId/nominate},
\texttt{POST\ /:submissionId/rerun} \\
detector \texttt{main.py} & 3 & \texttt{GET\ /health},
\texttt{POST\ /analyze}, \texttt{POST\ /features} \\
\end{longtable}

The web client presents 10 pages and 5 shared components. The detector
is reachable only from the API and only with a bearer token.

\hypertarget{non-functional-requirements}{%
\subsection{Non-functional
requirements}\label{non-functional-requirements}}

\begin{itemize}
\tightlist
\item
  \textbf{NFR-1 (security --- credential isolation)} No credential shall
  appear outside the API. The web client and the detector shall hold
  none.
\item
  \textbf{NFR-2 (security --- authorisation)} Every route shall enforce
  role and course-scope authorisation; a faculty member shall not read a
  submission outside their own courses.
\item
  \textbf{NFR-3 (performance --- bounded comparison)} The detector shall
  accept at most 50 candidates per request and shall run full tree
  comparison on at most 10 of them, selected by a cheap label-bag
  prefilter. This bounds the cost of one analysis irrespective of cohort
  size, and it is the behaviour any measurement of Layer 1 must
  reproduce to be faithful.
\item
  \textbf{NFR-4 (performance --- upload bound)} Submitted files shall be
  limited to the configured maximum (262,144 bytes as configured) and
  the detector request body to 16 MB.
\item
  \textbf{NFR-5 (reliability --- queue)} Detection shall be queued
  rather than synchronous, and debounced by job identity.
\item
  \textbf{NFR-6 (auditability)} A result shall be interpretable from
  stored data alone, without re-running detection and without reference
  to the current configuration.
\item
  \textbf{NFR-7 (evidence quality)} A reported match shall meet a
  minimum span size (10 nodes) before being shown to a reviewer, because
  two one-line functions resemble each other by virtue of being one line
  long.
\item
  \textbf{NFR-8 (explainability)} Every feature contributing to a
  behavioural score shall be explainable to a student contesting a flag.
  This constrains the feature set to a size a person can be walked
  through, and is the reason four features are scored rather than the
  ten extracted.
\item
  \textbf{NFR-9 (stated limitation --- Layer 2 efficacy)} Layer 2's
  ability to detect \textbf{partial} outsourcing is \textbf{not
  established}. On the 214-pair style-shift corpus (94 distinct source
  files) the layer did not separate spliced files from clean ones at any
  splice fraction, and combining it with Layer 1 degraded ranking
  precision. The cause is understood and structural: the four scored
  features are whole-file aggregates, and the score is a magnitude of
  deviation, so a partial substitution displaces a file in a direction
  that is as likely to move it toward its baseline as away. Any
  requirement asserting that Layer 2 detects outsourcing must be read
  against this. The design direction for addressing it --- per-unit
  rather than whole-file features, and a signed rather than magnitude
  comparison --- is future work and is not built.
\item
  \textbf{NFR-10 (stated limitation --- Layer 1 efficacy)} No precision
  or recall figure exists for structural similarity on
  \emph{known-copied} pairs, because no labelled corpus of genuine
  student plagiarism exists. The measured distribution is over honest
  work by different authors, which is a different claim.
\end{itemize}

\hypertarget{data-requirements}{%
\subsection{Data requirements}\label{data-requirements}}

Eight document collections.

\begin{longtable}[]{@{}
  >{\raggedright\arraybackslash}p{(\columnwidth - 4\tabcolsep) * \real{0.3333}}
  >{\raggedright\arraybackslash}p{(\columnwidth - 4\tabcolsep) * \real{0.3333}}
  >{\raggedright\arraybackslash}p{(\columnwidth - 4\tabcolsep) * \real{0.3333}}@{}}
\toprule\noalign{}
\begin{minipage}[b]{\linewidth}\raggedright
id
\end{minipage} & \begin{minipage}[b]{\linewidth}\raggedright
collection
\end{minipage} & \begin{minipage}[b]{\linewidth}\raggedright
holds
\end{minipage} \\
\midrule\noalign{}
\endhead
\bottomrule\noalign{}
\endlastfoot
\textbf{DR-1} & \texttt{User} & accounts and roles: admin, faculty,
student \\
\textbf{DR-2} & \texttt{Course} & a course and its enrolment \\
\textbf{DR-3} & \texttt{Assignment} & an assignment, its language, its
provenance \\
\textbf{DR-4} & \texttt{Submission} & an uploaded file reference, its
hash, its status \\
\textbf{DR-5} & \texttt{DetectionResult} & the score, both layers'
evidence, the applied weights and threshold, the recorded review \\
\textbf{DR-6} & \texttt{DetectionConfig} & per-course weights and
thresholds, with per-language overrides \\
\textbf{DR-7} & \texttt{BaselineProfile} & a student's style baseline,
per language \\
\textbf{DR-8} & \texttt{AuditLog} & every state change affecting a
student's record \\
\end{longtable}

\begin{itemize}
\tightlist
\item
  \textbf{DR-9} \texttt{DetectionResult} shall carry \texttt{isCurrent}
  and \texttt{revision}, so that a re-run supersedes without destroying.
\item
  \textbf{DR-10} \texttt{AuditLog} shall be independent of the documents
  it describes (see FR-32).
\end{itemize}

\begin{center}\rule{0.5\linewidth}{0.5pt}\end{center}

\hypertarget{verification}{%
\section{Verification}\label{verification}}

Requirements are verified by the automated suite described in
\texttt{docs/TESTING.md}: 155 API tests and 123 detector tests. The
suite is run in CI on every push, including the API integration tests,
which were added to the gate on 01 November 2026.

Chapter 6 of the report carries the test-case table mapping requirement
identifiers to tests, together with the mutation-testing results: 19
deliberate mutations across the detector suite, of which two were not
caught on the first attempt and exposed a dead code path and a missing
sort assertion.

\textbf{Requirements NFR-9 and NFR-10 are verified as limitations rather
than as capabilities}, by the measurement scripts in
\texttt{apps/detector/measurements/}, each of which is committed and
re-runnable.
